\chapter{Theory}
\label{chap:theory}

This chapter develops four theoretical foundations the rest of the thesis stands on. \Cref{sec:scheduling} sets out the scheduling problem the algorithm solves and names the utilization dimensions grouped under \textbf{Efficiency}. \Cref{sec:hitl} sets out the human-in-the-loop pattern under which the coordinator runs the system, and names \textbf{Trust/control} as a structural requirement of that pattern. \Cref{sec:tms} places RP inside the existing software category and identifies the planning gap RP fills. \Cref{sec:dsr} sets out the research paradigm that frames both the artefact and the knowledge claims made about it. Of the four, \Cref{sec:hitl} carries the most theoretical weight: the trust-and-control argument the rest of the thesis turns on is built there.

\section{Resource Scheduling}
\label{sec:scheduling}

Many everyday jobs come down to scheduling. A dentist fits one-to-one appointments with their patients into the hours of the working day. A waiter routes the next free server to the next ready table. A parent juggles pickup times across two children's after-school activities. Each is a small version of the same problem: matching a limited set of resources to tasks within a specific time window, under rules about who can do what. \textcite{pinedo2016scheduling} defines scheduling more formally as allocating resources to tasks over time, subject to given constraints, with the goal of optimising one or more objectives. The version that runs at a Norwegian transport company every morning sits inside this definition: a traffic coordinator matches drivers and vehicles to the day's transport assignments, where each assignment has a pickup time, a destination, and rules about who and what can do it.

The version of the problem RP solves is harder than the simplest single-resource case, because each transport assignment needs both a driver and a vehicle, and the two must stay paired for the whole assignment. In scheduling vocabulary, assignments are tasks with fixed time windows and pickup-and-delivery locations; drivers and vehicles are resources with their own attributes (a driver has competencies and availability, a vehicle has type and capacity). The planning horizon itself is set by the coordinator, from a single day to a week or longer. Pairing makes the problem harder than the single-resource case that dominates \parencite{pinedo2016scheduling}, because a candidate plan now has to satisfy compatibility for two resource types at the same moment.

Producing a plan that legally fits every assignment is one bar; producing one that uses the company's drivers and vehicles well, without unnecessary overtime, without long gaps between assignments, and without piling all the heavy work on a few drivers, is a different one, and harder. The personnel-scheduling literature judges plans not only by feasibility but by how well they use the resources available, on dimensions like overtime hours, idle time between assignments, and how evenly the workload spreads across the team \parencite{ernst2004staff}. These are not the only dimensions a transport company watches, but they are the ones the system most directly affects, and the thesis groups them under \textbf{Efficiency}. None of them can improve before the coordinator can see them: the system's first contribution is putting those quantities in front of the coordinator in a single view they can read at a glance.

The constraints that shape the plan split into the two kinds named in \Cref{sec:problem}, and the solver treats them differently. Hard constraints must hold for a plan to be feasible at all: the driver's licence class matches the vehicle, the driver holds any required certification such as the licence for transporting dangerous chemicals, the driver is on duty rather than on declared time-off, and no resource is double-booked. Soft constraints encode preferences the plan should honour where possible but may be traded off, such as even workload across drivers, individual driver preferences, and continuity with the same customer over time. Hard constraints bound feasibility; each violated soft constraint adds a penalty to the score the solver tries to minimise \parencite{ernst2004staff, rossi2006constraint}. The size of each penalty is a weight set per company in RP, so a coordinator's notion of a good plan can be expressed without code changes. That configurability is the artefact-level mechanism through which \textbf{Adaptability} is delivered, and \Cref{sec:iter-weights} describes how the weights are exposed.

Even with hard and soft constraints sharply defined, finding a guaranteed-best plan is out of reach in practice. The problem is NP-hard \parencite{rossi2006constraint, pinedo2016scheduling}, which means that no algorithm is known to solve every instance in time that grows polynomially with the input size, and one is widely believed not to exist. As the number of drivers, vehicles, and assignments grows, the time required to search every possible plan grows so fast that exhaustive search stops being practical even on the daily fleets the interviewed companies operate. RP responds by giving up provable optimality and using solvers that return good plans within a time budget the coordinator sets, and \Cref{sec:iter-multiengine} describes how three different engines were combined to do that.

When a problem is NP-hard, the practical question is no longer whether to find a guaranteed-best solution but which kind of algorithm to use to find a good one. Several broad styles exist. A teacher putting together a class timetable, who picks the first room that fits each lesson and moves on, is using one. A person solving a sudoku, who writes down what each cell could hold and rules out options as they go, is using another. A person editing a complicated calendar, who tries small changes and keeps the ones that improve the schedule, is using a third. The personnel-scheduling literature names these three styles formally. Constructive heuristics, called dispatching rules in \textcite{pinedo2016scheduling}, build a plan one assignment at a time according to a priority rule; they are fast, but can lock onto a locally good plan and miss a much better plan reachable only by reordering earlier choices. Constraint programming formalises the problem as variables and constraints and then searches the space of feasible plans systematically, using propagation to deduce forced choices and backtracking to undo dead ends; within a time budget it returns near-optimal plans, at the cost of higher modelling effort upfront \parencite{rossi2006constraint}. Metaheuristics improve a candidate plan by repeatedly modifying it and keeping the better candidate; the tabu search of \textcite{glover1986future} keeps a short-term memory of recent moves to avoid cycling, and the late-acceptance hill climbing of \textcite{burke2017late} accepts a new candidate if it is no worse than the candidate seen a fixed number of steps ago. Mathematical programming, a fourth body that dominates the rostering literature by volume, is rarely used on its own for messy real-world instances because their constraints and objectives resist clean mathematical formulation \parencite[pp.~17--18]{ernst2004staff}.

The three families differ along two axes that matter for RP: speed (how quickly they return a candidate plan) and quality (how close that plan is to optimal). Heuristics are fastest but can lock in early; constraint programming searches systematically but takes longer to set up; metaheuristics keep improving as long as time allows. RP runs all three behind one interface, and the coordinator chooses the speed-quality trade-off at the point of plan generation. \Cref{sec:iter-multiengine} reports how each engine behaves on the same benchmark instances. Solver capability alone, however, says nothing about who reads the resulting plan, who can change it, and on what authority. \Cref{sec:hitl} turns to that question.


\section{Human-in-the-Loop Automation}
\label{sec:hitl}

A nurse confirms a medication dose the hospital system flags. A pilot accepts an autopilot suggestion before the autopilot acts. In each case an automated process generates a recommendation, and a human exercises authority over whether it takes effect. The human factors literature calls this configuration human-in-the-loop. A traffic coordinator who reviews and approves every plan the algorithm produces is operating in exactly the same pattern.

\textcite{parasuraman2000automation} formalise this division of authority on a ten-level scale of automation for decision and action selection. The scale runs from level 1 (fully manual) to level 10 (fully autonomous), with intermediate steps that progressively transfer authority from human to machine. Level 4 is ``the computer suggests one alternative'', level 5 is ``executes that suggestion if the human approves'', and level 6 is ``executes automatically unless the human vetoes within a time window''. RP sits at level 5: the solver produces one plan, the plan is shown on the timeline, and no driver is dispatched until the coordinator approves it. The artefact deliberately stops short of level 6, where the coordinator would only need to fail to object.

As introduced in \Cref{sec:problem}, the reason for placing partial authority with the coordinator, rather than letting the algorithm act on its own, has a structural answer in \textcite{bainbridge1983ironies}. Bainbridge's central observation is that the people who specify automation are usually not the same as the people who must operate it: the owner pays for the system and asks for it, while the worker on the floor lives with it day to day. The split is concrete in the present project. Admmit's customers are the transport-company owners, who want planning automated to improve resource utilization. Admmit also speaks regularly with the coordinators those owners employ, and the coordinators want to keep control of every assignment made on their behalf. Both groups are real stakeholders for Admmit, and a system that satisfied only one would not be acceptable. The human-in-the-loop operating model is the design move that satisfies both at once: the owners get the automation that improves utilization, and the coordinators get the review-and-override authority that keeps them in charge. \textbf{Trust/control} is therefore a structural requirement of the system, not a usability feature added on top, and it follows directly from who is asking for the system and who has to live with its mistakes.

Trust in the algorithm is a separate question from authority over it. A coordinator who has the authority to override the algorithm may still leave the system unused if they do not trust what it produces. \textcite{hoff2015trust} review 127 studies of human-automation trust from 101 articles, and group their findings into a three-layer model of dispositional, situational, and learned trust \parencite[p.~413]{hoff2015trust}. Dispositional trust reflects how much an individual tends to trust automation in general; situational trust depends on context and workload; learned trust is built through experience with the specific system, first as initial trust based on reputation, then as dynamic trust that updates with each interaction \parencite[p.~420]{hoff2015trust}. Two coordinators using the same plan generator may therefore hold different levels of trust at the same moment, and the same coordinator's reliance on the algorithm shifts across days. Hoff and Bashir also report that early errors damage trust more than later errors do \parencite[p.~426]{hoff2015trust}, which has a concrete consequence for deployment: a coordinator's first week with RP weighs more in trust formation than the months that follow. The aim is calibrated trust, where the coordinator neither rubber-stamps the algorithm nor refuses to use it.

For the coordinator to trust the algorithm in a way that matches how good it actually is, they need to understand what it is doing. \textcite{miller2019explanation} surveys explanation in the social sciences and notes three things most automated systems get wrong. First, explanations are contrastive: when people ask why something happened, they usually mean why it happened rather than the alternative they had in mind. Second, explanations are selective: people want one or two key reasons, not the full causal chain. Third, explanations are social: explaining is a kind of communication, so a good explanation fits what the listener already knows \parencite[p.~3]{miller2019explanation}. Translated to a planning algorithm, the coordinator's question is rarely ``why was this assignment generated?''. It is ``why was this driver assigned, rather than the one I would have picked?''. Miller argues that explanation belongs in the interface itself, and that ``trust is lost when users cannot understand traces of observed behaviour or decisions'' \parencite[p.~3]{miller2019explanation}. The implication for RP is that the system has to give reasons in the coordinator's vocabulary (competence, time, current load), not in raw constraint weights and solver internals, and those reasons belong next to the override controls themselves.

\textcite{lee2004trust} define trust as the user's belief that the system will help them reach their goals in a situation where they cannot fully predict what it will do \parencite[p.~54]{lee2004trust}. Their key argument is that the design goal is the right amount of trust, not the maximum amount. A coordinator who accepts every plan without checking it (Lee and See call this \textit{misuse}) is over-trusting; a coordinator who plans entirely by hand even when the solver works well (\textit{disuse}) is under-trusting. To help a user calibrate, Lee and See recommend that the system reveal its past performance, show how it works, and explain its purpose in terms that connect to the user's goals \parencite[p.~74]{lee2004trust}. The risk Lee and See point to is concrete: if coordinators do not understand or trust the algorithm, they will not use it, and the visibility gap from \Cref{sec:scheduling} stays open regardless of how good the solver is. \Cref{chap:discussion} returns to this adoption question.

The four bodies above (Parasuraman's automation taxonomy, Hoff and Bashir's trust antecedents, Miller's explanation as interface, and Lee and See's definition of appropriate trust) translate into concrete design constraints on RP's HITL interface. The planning timeline shows the algorithm's suggested plan in the same Gantt view the coordinator uses for manual work, so the suggestion is open to inspection rather than presented as a finished output. After a solver run, a summary panel reports the solution score and counts of hard and soft violations, exposing process and performance in the sense Lee and See describe. A separate deviation view lists each active conflict by type, affected assignment, and affected resource, so a coordinator who disagrees with the algorithm can locate the contested assignment and decide whether to override. Override is intentionally low-friction: dragging an assignment block to a different employee row reassigns it, and conflicting changes save with a warning rather than a block. A blocking validator would tell the coordinator that the system, not the coordinator, holds final authority, which is the configuration Bainbridge warns against and the one \textbf{Trust/control} forbids. Soft-constraint weights are exposed per company so that a coordinator's notion of a good plan can be adjusted without code changes, and the choice of solver engine and time budget is exposed at the point of plan generation so the system's capability limits are visible before the plan is read. \Cref{sec:iter-hitl,sec:iter-deviations,sec:iter-tradeoff} document how each layer was built into the artefact.


\section{Transport Management Systems}
\label{sec:tms}

Transport companies run their day-to-day operations on a category the industry groups under one label: a Transport Management System (TMS). \textcite{griffis2007tms} define a TMS as information technology used to plan, optimise, and execute transportation operations before, during, and after the physical movement of goods \parencite[p.~19]{griffis2007tms}. \textcite{heinbach2022datadriven} adopt a compatible definition, traceable in their work to TMS as freight-forwarder software for assigning customer shipments to fleet resources \parencite[p.~811]{heinbach2022datadriven}.

Inside that label, the picture is not uniform. What a TMS actually offers depends on the vendor, and whether a transport company uses one at all depends on the company. In the Norwegian market, the two commercial systems named in the interviews are \textbf{Timpex} and \textbf{Opter}; both cover order management and freight billing, and neither generates assignment plans automatically. Adoption itself is uneven: \textcite{griffis2007tms} report that of the firms in their 2007 North American sample without any TMS, half manage transportation entirely manually, while the other half rely on partial or legacy IT \parencite[p.~26]{griffis2007tms}. The sector RP enters is therefore a patchwork, with different vendors covering different ground and a substantial share of operations running without dedicated planning software at all.

Even where a TMS is in place, one part of the operation is missing from the category as a whole: the upstream planning step that decides which driver and vehicle perform which assignment. \textcite{griffis2007tms} observe that most TMS offerings centre on the individual shipment as the unit of analysis and lack the ability to optimise multi-load shipments, leaving load consolidation as a manual activity \parencite[p.~27]{griffis2007tms}. \textcite{heinbach2022datadriven} reach the same conclusion: TMS products focus on shipment-level business activities, while other digital platforms add capabilities outside that focus \parencite[p.~808]{heinbach2022datadriven}. What existing systems cover, in short, is order management, freight billing, and shipment-level execution; what is missing is the upstream planning step between an order existing and a specific driver and vehicle being assigned to perform it. RP is positioned in exactly that gap. \Cref{chap:findings} substantiates the gap with the empirical comparison, and the per-company configurability that lets RP fit different vendor environments is built in \Cref{sec:iter-weights} and discussed under \textbf{Adaptability} in §5.1.3. \Cref{sec:dsr} turns to the research paradigm under which an artefact filling that gap is built and investigated.


\section{Design Science Research}
\label{sec:dsr}

Design Science Research (DSR) is the research paradigm in which a designed artefact is constructed and investigated as the central contribution of a study, together with the knowledge produced about how that artefact behaves in its context. The point is that the artefact and the knowledge come out of the same project. \textcite{wieringa2014dsm} states the position concisely: design science is the design and investigation of artefacts in context, where the artefact is built to improve something in its problem context \parencite[p.~3]{wieringa2014dsm}.

\textcite{hevner2007threecycle} structures the same work as three linked cycles: the relevance cycle imports requirements and acceptance criteria from the application environment, the rigor cycle draws on existing scientific knowledge and contributes new knowledge back to it, and the design cycle iterates internally between constructing and evaluating the artefact \parencite[pp.~88--90]{hevner2007threecycle}. \textcite{peffers2007dsrm} structure the same paradigm as a procedural model. Their Design Science Research Methodology (DSRM) specifies six activities: problem identification and motivation; objectives for a solution; design and development; demonstration; evaluation; and communication \parencite[p.~46]{peffers2007dsrm}. The two framings are complementary: the cycles describe what must be present, the activities describe what is done. \Cref{chap:method} applies the DSRM activities step by step to this project.

DSR is not the only research paradigm available for an information systems study. \textcite{orlikowski1991studying} group IS research into three philosophical stances: positivist (assuming an objective social reality measurable with structured instruments), interpretive (treating organisational reality as socially constructed and reading meaning through participants' frames), and critical \parencite[p.~5]{orlikowski1991studying}. A purely positivist study fits population-level measurement of a deployed system, but cannot on its own justify building a new artefact whose primary contribution is the design knowledge embedded in it. A purely interpretive study fits rich description of coordinator practice, but cannot on its own check a runnable artefact's behaviour against external criteria. DSR sits between the two: qualitative input informs requirements, and the artefact itself is examined against those requirements.

Inside DSR, the distinction between validation and evaluation shapes what kind of empirical evidence the project can offer. \textcite{wieringa2014dsm} distinguishes validation, which predicts whether a deployed system would do what stakeholders want, from evaluation, which observes a deployed system in actual use \parencite[pp.~31, 34]{wieringa2014dsm}. In validation, the researcher experiments with a model of how stakeholders would use the artefact; in evaluation, stakeholders use the artefact independently of the researcher. The present thesis validates rather than evaluates RP. Behaviour is predicted in two ways: benchmarking on synthetic datasets, which is a single-case mechanism experiment in Wieringa's classification (one artefact, controlled inputs) \parencite[p.~64]{wieringa2014dsm}, and requirements traceability against the requirements derived in the relevance cycle. No transport company has operated the system independently. \textcite{hevner2007threecycle} treats this ordering as expected practice (artefacts tested in laboratory and experimental situations before being released into field testing along the relevance cycle) \parencite[p.~91]{hevner2007threecycle}, so the absence of production deployment is a structural feature of the project's stage rather than a methodological deficit. \Cref{sec:evaluation-framework} applies this distinction to the synthetic-benchmark protocol, and Section 5.4 returns to it as limitation L5 when bounding what that validation can and cannot claim.
